% docs/math/ensemble_adjoint.tex
% The particle-carried ensemble adjoint of JAX-DSMC, as implemented on branch
% fix/paper-adjoints (2026-09-17).  Build: see docs/math/README.md.
%
% Equation numbers are section-scoped (e.g. (3.7)) and are meant to be cited from
% code comments as "docs/math/ensemble_adjoint.tex eq. (3.7)".  Do not renumber
% sections without updating those comments.
\documentclass[11pt,a4paper]{article}
\usepackage{amsmath}
\usepackage{amssymb}
\usepackage[colorlinks=true,linkcolor=blue,urlcolor=blue,citecolor=blue]{hyperref}

\numberwithin{equation}{section}

% page layout without the geometry package
\setlength{\textwidth}{16cm}
\setlength{\oddsidemargin}{0cm}
\setlength{\evensidemargin}{0cm}
\setlength{\textheight}{24cm}
\setlength{\topmargin}{-1cm}

\newcommand{\E}{\mathbb{E}}
\newcommand{\R}{\mathbb{R}}
\newcommand{\ind}{\mathbf{1}}
\newcommand{\code}[1]{\texttt{#1}}
\newcommand{\pth}{\mathrm{post}}
\newcommand{\pre}{\mathrm{pre}}
\newcommand{\Lam}{\Lambda}
\newcommand{\dd}{\mathrm{d}}
\newcommand{\Nn}{\mathcal{N}}
\newcommand{\pref}[1]{P3~eq.~\code{#1}}

\title{The ensemble adjoint of JAX-DSMC as implemented\\
\large branch \code{fix/paper-adjoints}, 2026-09-17}
\author{}
\date{}

\begin{document}
\maketitle

\noindent
This note states, equation by equation, what the particle-carried ensemble adjoint
in \code{jax\_dsmc/ensemble\_backward.py} (host reference) and
\code{jax\_dsmc/ensemble\_device.py} (JIT port) computes, where each term comes from
in the reference paper (``P3'', \code{papers\_reference/arXiv-2603.20946v2/main.tex},
section~5, \emph{Ensemble-Adjoint DSMC Method With Inflow Boundary Conditions}), and
where in the code it lives.  Equation labels of the form \pref{label} are the
\verb|\label{eq:label}| names in P3's source.  Numbers quoted in section~\ref{sec:validation}
are copied from \code{docs/ADJOINT\_BRANCH\_REVIEW.md} (``review'') sections~5.7, 5.12 and~6.
Everything marked \textbf{planned} is not in the code.

\section{Forward step and notation}
\label{sec:forward}

One forward step $k \to k+1$ of \code{simulate.\_step} (\code{jax\_dsmc/simulate.py})
is, in order,
\begin{equation}
\label{eq:step-order}
(x_k, v_k, a_k)
\;\xrightarrow{\text{move + walls}}\;
(x^m_k, v^m_k, a^m_k)
\;\xrightarrow{\text{inject}}\;
(x^m_k, v^m_k, a^m_k) \cup \{\text{newborns}\}
\;\xrightarrow{\text{sort}}\;
\text{cell tables}
\;\xrightarrow{\text{collide}}\;
(x_{k+1}, v_{k+1}, a_{k+1}),
\end{equation}
with $x_{k+1} = x^m_k$ (collisions never move particles) and $a$ the boolean
alive mask over a fixed number $N$ of particle slots.  This is the reverse of P3's
order (Algorithm~1 there: collide, then \pref{x\_update} transport), which flips
the backward order in section~\ref{sec:adjoint}.
The superscript $m$ denotes the \emph{post-move, post-injection, pre-collision}
state: it is what \code{EventRecord.x}, \code{.v}, \code{.alive} store
(\code{simulate.EventRecord}), while \code{EventRecord.v\_post} $= v_{k+1}$.

\paragraph{Move and walls.}
Free flight $x^m = x + \Delta t\, v$ with chronological first-event tracking on every
face (\code{wall\_adjoint.\_advance\_impl}); specular and periodic faces are
deterministic maps; an absorbing/outflow face sets $a^m_i = 0$; a diffuse face
re-emits with the accommodation-mixed kernel of section~\ref{sec:scores}.
The flight budget is the $\varepsilon$-regularised time $\tau_i = \max(\Delta t +
\varepsilon z_i, 0)$, $z_i \sim \mathcal N(0,1)$ (P3 section~4's random-time
wall, \pref{x\_update\_new}); $\varepsilon = 0$ recovers the deterministic step.

\paragraph{Inject.}
For every \code{INFLOW} face $f$ the expected count is (\code{inflow.target\_count})
\begin{equation}
\label{eq:target-count}
\bar N_f(\theta) = n_{\rho,f}\, v_{s,f}\, F(s_f)\, A_f\, \Delta t / F_{\mathrm{num}},
\qquad
v_{s,f} = \sqrt{2 k_B T_f/m},\quad
s_f = \frac{u_f\cdot n_f}{v_{s,f}},
\end{equation}
where $F$ is the untruncated flux integral (\code{inflow.flux\_integral}) and $n_f$
the inward normal.  The realised count is the stochastic rounding
$N_f = \lfloor \bar N_f \rfloor + \ind[u < \operatorname{frac}\bar N_f]$
(\code{sampling.stochastic\_round}).  Each newborn draws a velocity $v_b$ from the
normalised inward flux density $b_f$ of~\eqref{eq:flux-density}, a face position, and
a first-move fraction $\xi \sim \mathcal U[0,\Delta t]$, and is placed at
$x_f + \xi v_b$ with the same wall tracking as everybody else
(\code{inflow.inject}, \code{track=True}).  This is \pref{ensemble\_inflow\_left\_state}
with P3's $g_{k,i}$ replaced by the drifting half-Maxwellian sample.  The record
keeps $v_b$ \emph{as sampled} (\code{born\_v}, the score's argument), the face
(\code{born\_face}) and the newborn's own birth-step wall hits (\code{born\_walls}).

\paragraph{Collide (NTC, VSS).}
Per cell $c$ with $N_c$ members and volume $V_c$, the attempt count is the
stochastic rounding of
\begin{equation}
\label{eq:attempts}
\bar n_{\mathrm{att},c} = \tfrac12 N_c (N_c - 1)\, \Sigma_c\, \Delta t\, F_{\mathrm{num}} / V_c
\qquad (\code{ntc.attempt\_count}),
\end{equation}
where $\Sigma_c$ is the per-cell majorant (thermal a-priori or certified,
\code{ntc.thermal\_majorant}, \code{ntc.certified\_majorant}), recorded as
\code{EventRecord.sigma}.  Each attempt $a$ of cell $c$ (\code{ntc.\_slot}, one scan
slot for all cells) draws an unordered pair $(i,j)$ uniformly, accepts it with probability
\begin{equation}
\label{eq:accept}
p_{ij} = \min\!\big(q(g_{ij})/\Sigma_c,\, 1\big),
\qquad
q(g) = \sigma_T(g)\, g = C_{\mathrm{VSS}}\, (g^2)^{1-\omega},\quad g = |v_i - v_j|
\qquad (\code{vss.collision\_rate}),
\end{equation}
and on acceptance replaces $(v_i, v_j)$ by the VSS post-collision pair
$(v_i', v_j') = \operatorname{scatter}(v_i, v_j; u_\chi, u_\epsilon)$
(\code{vss.scatter}, deflection law in $\alpha$).  Attempts are processed in
slot-index order, so a particle may collide more than once per step and the
$a$-th attempt sees the velocities left by attempts $< a$.  The record stores, per
slot and cell, $(i, j, \text{accept}, v_i^{\pre}, v_j^{\pre}, v_i^{\pth}, v_j^{\pth})$
(\code{ntc.PairRecord}).  The accepted pairs are exact samples of the collision
measure $q(g) f f$ if and only if the majorant is never violated
(\code{Diagnostics.majorant\_violations} $= 0$).

\paragraph{RNG contract.}
Every stage of every step draws from
\begin{equation}
\label{eq:key-for}
\code{key\_for}(\kappa, k, s) = \operatorname{fold\_in}\big(\operatorname{fold\_in}(\kappa, k), s\big),
\qquad s \in \{\code{MOVE}, \code{WALL}, \code{NTC\_ROUND}, \code{PAIR}, \code{ACCEPT}, \code{SCATTER}, \code{INFLOW}\},
\end{equation}
(\code{jax\_dsmc/rng.py}), with the attempt-slot index folded in once more for
per-slot streams.  The capture (\code{simulate.run(capture=True)}) does not change any
stream or any arithmetic of the forward
(\code{tests/test\_ensemble\_backward.py::test\_capture\_is\_consistent\_with\_the\_states}).
The backward's own randomness (virtual partners, section~\ref{sec:virtual}) uses a
separate key.

\section{Objectives}
\label{sec:objective}

\paragraph{Terminal objective (implemented).}
\begin{equation}
\label{eq:objective}
J = \frac{1}{\Nn} \sum_{i:\, a_{M,i} = 1} r(x_{M,i}, v_{M,i}),
\end{equation}
which is \pref{ensemble\_inflow\_objective} with one difference: P3 fixes
$\Nn = N_0 + M(N_L + N_R)$; here $\Nn$ is the \code{normalization} argument of
\code{ensemble\_backward.backward} and is a constant of the run (the nominal particle
number $n_{\mathrm{nominal}}$ on the inlet shock, $N$ on Couette and the relaxation).
The per-particle reward $r$ is the callable \code{phi(x, v)} on the host; the device
takes the already-divided vector \code{phi\_final} $= r/\Nn$.

\paragraph{Value form.}
The device and the host both accept an arbitrary per-particle terminal vector, which
is used for non-additive objectives by the delta method: for the density ratio
$J = N_{\mathrm{down}}/N_{\mathrm{up}}$ (\code{examples/ensemble\_inlet\_shock.py::objective\_terminal})
\begin{equation}
\label{eq:delta-method}
\lambda_{M,i} = \frac{\ind[x_{M,i} \in \text{down}]}{N_{\mathrm{up}}}
 - \frac{N_{\mathrm{down}}}{N_{\mathrm{up}}^2}\,\ind[x_{M,i} \in \text{up}],
\end{equation}
with the two totals frozen at their realised values ($O(1/N)$ bias, review section~6).

\paragraph{Windowed and per-hit rewards (value form, implemented 2026-09-18).}
For $J = \Nn^{-1}\sum_{k} \sum_i \phi_k(x^{\pth}_{k,i}, v^{\pth}_{k,i})
 + \Nn^{-1}\sum_k \sum_{\text{hits } j \text{ of step } k} r_{k,j}$
the ensemble adjoint acquires a running source, in the sense of P3 section~7.5
(\pref{running\_alpha}--\pref{running\_beta} are its pathwise analogue).
\code{backward} and \code{backward\_virtual} take \code{phi\_step} $=\phi_k/\Nn$
(\code{f64[M, N]}, at each step's post-collision state -- the state
\code{simulate.run}'s \code{reward} hook sees) and \code{hit\_reward} $=r_{k,j}/\Nn$
(\code{f64[M, K, N]}, one slot per tracker pass).  Before the pull-back of step $k$
\begin{equation}
\label{eq:running-source}
\mu_{k,i} \leftarrow \mu_{k,i} + \phi_k(x^{\pth}_{k,i}, v^{\pth}_{k,i})/\Nn ,
\end{equation}
and on the move pull-back of step $k$ (survivors only; a newborn's rewards go to its
source score)
\begin{equation}
\label{eq:hit-source}
\mu_{k-1,i} = \lambda^m_{k,i} + \frac{1}{\Nn}\sum_{j} r_{k,j,i} .
\end{equation}
The weight of the wall-kernel score~\eqref{eq:wall-score} of hit $j$ is the
reward-to-go \emph{from that draw onward},
\begin{equation}
\label{eq:hit-weight}
w_{k,j,i} = \lambda^m_{k,i} + \frac{1}{\Nn}\sum_{j' \ge j} r_{k,j',i},
\end{equation}
the hit's own reward included -- it is a function of the drawn velocity (leaving it
out made the Couette windowed-shear gradient with respect to the wall speed read
$-0.0004$ against an FD of $-0.0165$; with it, $-0.0177 \pm 0.0014$).  The terminal
vector \code{phi\_final} is unchanged and equals $\phi_M/\Nn$ for a windowed objective
whose last step is the final one.  Gates: G1 (windowed $v_y^2$, \code{speed}, $z=-0.05$)
and G2 (windowed shear as a per-hit reward, \code{speed} $z=-0.77$, $T_{\mathrm{left}}$
$z=+0.46$), 48 seeds each against CRN finite differences
(\code{examples.ensemble\_couette --objective}).

\section{The ensemble adjoint as implemented}
\label{sec:adjoint}

\subsection{Carried values and the terminal condition}

P3 defines $\lambda_k(x,v) = \delta J/\delta f_k(x,v)$ with the terminal condition
\pref{ensemble\_inflow\_terminal\_adjoint}, $\lambda_M = r/\Nn$ in $\Omega$ and $0$
outside.  The code represents $\lambda_k$ by \emph{one scalar per particle slot at
that slot's own state} (route (b) of review section~5.1; P3 section~5.4 uses retained
subsets and reconstructs the field at every retained state instead):
\begin{equation}
\label{eq:terminal}
\mu_i := \lambda_{M}(x_{M,i}, v_{M,i}) = a_{M,i}\; r(x_{M,i}, v_{M,i})/\Nn
\qquad
\text{(\code{backward}: \code{mu = where(alive\_final, phi/normalization, 0)})}.
\end{equation}
Going backward through step $k$, the sweep turns $\mu$ (values at $t_{k+1}$, i.e.\ at
$(x^m_k, v_{k+1})$) into $\lambda^m_i := \lambda(x^m_{k,i}, v^m_{k,i})$ at the post-move
state, records what sources and walls need there, and then into $\mu$ at $t_k$.
The device port emits exactly $\lambda^m$ per step as \code{Backward.lam[M, N]}.

\subsection{Move pull-back}
\label{sec:move}

P3's transport pull-back \pref{ensemble\_inflow\_transport\_adjoint} is
$\tilde\lambda_{k+1}(x,v) = \ind[x + \Delta t\, v \in \Omega]\,\lambda_{k+1}(x + \Delta t\, v, v)$.
Because the value is carried on the particle that made the move, this is the
identity on survivors:
\begin{equation}
\label{eq:move-pullback}
\mu_i^{(k)} = \lambda_k(x_{k,i}, v_{k,i}) =
\begin{cases}
\lambda^m_i & a^m_{k,i} = 1 \text{ and slot } i \text{ was not born in step } k,\\
0 & \text{otherwise (absorbed during the move, or born this step)},
\end{cases}
\end{equation}
(\code{backward}: \code{mu = where(alive \& \textasciitilde born, lam, 0)};
device \code{\_step\_pullback}, \code{\_step\_pullback\_virtual}).
The identity holds because specular and periodic faces are deterministic maps
(the value moves with the particle) and a diffuse face is a Markov kernel whose
realised re-emission is an unbiased sample of it, so
$\E[\lambda(\text{after hit})\mid \text{before hit}]$ is estimated by the one carried
value (review section~5.6).  No $\varepsilon$ enters: the regularised flight time is
parameter-free randomness to this estimator for every control in
section~\ref{sec:validation}.

\subsection{Collision pull-back}
\label{sec:collision}

P3 pulls the adjoint through the collision step with the adjoint of the linearised
operator, \pref{ensemble\_collision\_adjoint\_inflow},
\begin{equation}
\label{eq:Lstar}
(\mathcal L^*_{f} \psi)(v) = \int_{\R^3} f(w) \int_{S^2} q(\sigma)
\big[\psi(v') + \psi(w') - \psi(v) - \psi(w)\big]\, \dd\sigma\, \dd w,
\end{equation}
and the one-step recursion \pref{ensemble\_adjoint\_recursion\_inflow},
$\lambda_k = \tilde\lambda_{k+1} + \Delta t\, \mathcal L^*_{f_k}\tilde\lambda_{k+1}$.
Because collisions here come last in the step, the code applies the collision
pull-back to $\mu$ (the values at $t_{k+1}$) \emph{before} the move pull-back of
section~\ref{sec:move}; the object it produces is $\lambda^m$ at the post-move state.
In the code $q$ is the VSS rate of~\eqref{eq:accept} and the angular measure is the
VSS deflection law, not P3's Maxwell kernel; with realised or NTC-weighted pairs the
kernel never has to be integrated explicitly.

Three pull-backs are implemented and selected process-wide by
\code{ensemble\_backward.configure(pullback, S, K, nodes)}.  Throughout,
$\Lam^{(\cdot)}$ denotes the per-cell field interpolant of section~\ref{sec:field}
built from the current carried values on the current node velocities.

\subsubsection{Sequential (exact discrete adjoint of the NTC step; reference)}
\label{sec:sequential}

Walk the accepted attempts of each cell in \emph{reverse} slot order.  At attempt
$a$ with pair $(i,j)$, pre-collision velocities $(v^{\pre}_i, v^{\pre}_j)$ from the record,
current carried values $\lambda_i, \lambda_j$ (which at that point are the value
function right after attempt $a$, at the post-collision states) and current node
velocities $V^{(a)}$ (the population as it stood right after attempt $a$):
\begin{equation}
\label{eq:sequential}
\begin{aligned}
\lambda_i &\leftarrow \lambda_i + \lambda_j - \Lam^{(a)}(x_j, v^{\pre}_j), \\
\lambda_j &\leftarrow \lambda_j + \lambda_i - \Lam^{(a)}(x_i, v^{\pre}_i), \\
\Phi_a &:= \lambda_i + \lambda_j - \Lam^{(a)}(x_i, v^{\pre}_i) - \Lam^{(a)}(x_j, v^{\pre}_j),
\end{aligned}
\qquad
\text{then } V^{(a-1)}_i \leftarrow v^{\pre}_i,\; V^{(a-1)}_j \leftarrow v^{\pre}_j,
\end{equation}
where the right-hand sides use the values before the update
(\code{ensemble\_backward.\_collision\_pullback\_sequential}; device
\code{ensemble\_device.\_slot\_events} inside a reverse \code{lax.scan} over slots).
This is the four-term sum of~\eqref{eq:Lstar} with the realised pair as the single
Monte Carlo sample of the collision measure and the two ``$+$'' terms read from the
carried values (no interpolation), so the only interpolated quantity per accepted
attempt is the partner's pre-collision ``ghost'' (review section~5.1).  With
\code{GHOST\_FIELD = "current"} the ghost is read from the field as chained so far,
which makes the sweep the exact discrete adjoint of the sequential NTC step
(\code{test\_sequential\_pullback\_reproduces\_the\_pre\_collision\_field\_for\_an\_invariant}).
$\Phi_a$ is the pair's \emph{four-term} (\code{Pairs.four\_term}, device \code{Backward.four\_term}), weight~1 per realised pair.
Its per-seed noise grows with the collisions per particle per step and does not fall
with $\Delta t$ (review section~5.10).

\subsubsection{\texorpdfstring{$\Delta t$-linearised (reference only, biased $O(\Delta t)$)}{dt-linearised (reference only, biased O(dt))}}
\label{sec:linearised}

All accepted attempts at once, every ghost from the end-of-step field $\Lam_{k+1}$
fitted on $(x^m, v_{k+1})$ with values $\mu$:
\begin{equation}
\label{eq:linearised}
\lambda^m_i = \mu_i + \sum_{a \ni i} \Big[\Lam_{k+1}(x_{j_a}, v^{\pth}_{j_a}) - \Lam_{k+1}(x_{j_a}, v^{\pre}_{j_a})\Big],
\qquad
\Lam_{k+1}(x_{j}, v^{\pth}_{j}) = \mu_j \text{ when } a \text{ was } j\text{'s last accepted attempt of the step},
\end{equation}
(\code{ensemble\_backward.\_collision\_pullback}, \code{configure("linearised")}).
This is \pref{ensemble\_adjoint\_recursion\_inflow} quadratured on the realised
pairs, and it is what P3's single-step recursion literally says; it is
$-8.6\%$ biased on the relaxation arbiter at $0.16$ collisions/particle/step
(review section~5.12) because a third of the accepted pairs have a member colliding
again in the same step and the end-of-step field is not the value after
\emph{their} attempt.  Not shipped.

\subsubsection{Virtual partners, substepped (shipped: \code{VIRTUAL = (S, K) = (4, 8)}, \code{VIRTUAL\_NODES = "staged"})}
\label{sec:virtual}

The backward collision operator $e^{\Delta t \mathcal L^*}$ is applied in $S$ substeps
of $\Delta t / S$ each, going backward $s = S-1, \dots, 0$.  Substep $s$ uses node
velocities $V^{(s+1)}$ (below) and one smooth per-cell fit $\Lam^{(s)}$ of the current
$\lambda$ on those nodes.  For every alive particle $i$ in cell $c$ and each of
$K$ virtual partners $j_\kappa$ drawn uniformly among the \emph{other} alive members of
$c$, with fresh VSS uniforms $(u_\chi, u_\epsilon)_{\kappa}$:
\begin{equation}
\label{eq:virtual}
\lambda_i \leftarrow \lambda_i + \sum_{\kappa=1}^{K} w_{i\kappa}
\Big[\Lam^{(s)}(x_i, v_i') + \Lam^{(s)}(x_{j_\kappa}, v_{j_\kappa}') - \Lam^{(s)}(x_i, v_i) - \Lam^{(s)}(x_{j_\kappa}, v_{j_\kappa})\Big],
\end{equation}
\begin{equation}
\label{eq:virtual-weight}
w_{i\kappa} = \frac{q(|v_i - v_{j_\kappa}|)}{\Sigma_c}\cdot\frac{2\, n_{\mathrm{att},c}}{N_c}\cdot\frac{1}{K S},
\qquad (v_i', v_{j_\kappa}') = \operatorname{scatter}(v_i, v_{j_\kappa}; u_\chi, u_\epsilon)_\kappa ,
\end{equation}
where $v_i = V^{(s+1)}_i$, $n_{\mathrm{att},c}$ is the realised attempt count of the cell
in this step (\code{pairs.pi < N}, summed over slots) and $N_c$ its alive occupancy
(\code{ensemble\_backward.\_collision\_pullback\_virtual};
device \code{ensemble\_device.\_substep}, \code{\_partner\_pass}).  All four values come
from the same fit, so their interpolation errors cancel to leading order.  This is
P3's Monte Carlo collision correction \pref{MC\_collision\_adjoint\_inflow},
$\widehat C_{k,i} = \Delta t\, \rho^j_k Q_0\, N_{\mathrm{col}}^{-1} \sum_s [\cdots]$, with
$N_{\mathrm{col}} = K$ and the Maxwell prefactor $\Delta t\, \rho Q_0$ replaced by the NTC
expectation weight $(q/\Sigma_c)\,(2 n_{\mathrm{att},c}/N_c)$: the expected number of
collisions a particle suffers in the step is $2 n_{\mathrm{att},c}/N_c$ attempts times
the acceptance $q/\Sigma_c$, and the realised attempt count already carries
$\Delta t$, $F_{\mathrm{num}}$, $V_c$ and $\Sigma_c$ from~\eqref{eq:attempts}.  No realised
pair enters except through $n_{\mathrm{att},c}$.  The substep count $S$ is not in P3;
it converts the single-step recursion into an $O(\Delta t/S)$ splitting
(review section~5.12: $+1.8\%$ at $S=2$, $+0.7\%$ at $S=4$, $+0.1\%$ at $S=8$ on the relaxation arbiter).

\paragraph{Staged nodes.}
The NTC step processes slots in index order, so the population at the boundary
after a fraction $b/S$ of a cell's attempts is $v^{\pre}$ overwritten, for each
particle, by the post-collision velocity of its last accepted slot before that
boundary (\code{ensemble\_backward.\_staged\_velocities}, device
\code{ensemble\_device.staged\_velocities}; tested against a sequential replay).
The nodes of substep $s$ are $V^{(s+1)}$, the state at the end of its own slot range,
with $V^{(0)} = v^m$ and $V^{(S)} = v_{k+1}$.  After the update~\eqref{eq:virtual} a
particle whose velocity differs across the boundary moves its value through a
\emph{refitted} field $\Lam'^{(s)}$ (current $\lambda$, nodes $V^{(s+1)}$):
\begin{equation}
\label{eq:boundary-move}
\lambda_i \leftarrow \lambda_i + \Lam'^{(s)}(x_i, V^{(s)}_i) - \Lam'^{(s)}(x_i, V^{(s+1)}_i)
\qquad \text{if } V^{(s)}_i \neq V^{(s+1)}_i
\end{equation}
(\code{\_boundary\_move} on the device).  With \code{nodes = "post"} every $V^{(s)}$,
$s \ge 1$, equals $v_{k+1}$ and the field inside the step is frozen ($+5$--$6\%$
biased, review section~5.12); the runners pass \code{"staged"} (note that
\code{backward\_virtual}'s own default argument is \code{"post"}).

\paragraph{Kernel-parameter term.}
Each virtual pair also returns its four-term $\Phi_{i\kappa}$ (the bracket
of~\eqref{eq:virtual}) with weight $\tfrac12 w_{i\kappa}$ (\code{Pairs.weight}; the
$\tfrac12$ because every unordered pair is generated from both members), and the
device accumulates per step
\begin{equation}
\label{eq:kernel-term}
\code{Backward.kernel}_k = \sum_{i,\kappa} \tfrac12\, w_{i\kappa}\, \Phi_{i\kappa}\; \partial_\theta \log q(v_i, v_{j_\kappa})
\end{equation}
when a \code{pair\_score} $= \partial_\theta \log q$ is supplied (used by
\code{examples/ensemble\_relaxation.py} for \code{diam} and \code{omega}).

\paragraph{Stability.}
The quadrature noise injected by~\eqref{eq:virtual} is smoothed by the next fit; with
too few partners for the cell occupancy the carried field grows geometrically over a
long horizon ($K=2$ at 40 particles per cell blew up by $3\times 10^6$ over 300 steps;
$K \ge 8$ decays, review section~5.12).  \code{examples.ensemble\_inlet\_shock.device\_backward}
refuses a backward whose rms $\lambda$ grew more than \code{GROWTH\_LIMIT} $= 100\times$
over the horizon.

\subsection{The field interpolant}
\label{sec:field}

P3 uses a weighted local \emph{affine} fit in $(x, v_1, v_2, v_3)$ on retained
particles (section~5.4 there).  Here the interpolant is cell-local and quadratic.
For a query $(x_q, v_q)$ in cell $c$, the nodes are the alive particles of $c$ with
the current node velocities, in the scaled coordinates
$\tilde p = (x/\Delta x,\; v/v_{\mathrm{scale}})$ (cell width per axis; a thermal speed).
Take the $k = 32$ nearest nodes in that metric (\code{neighbors=32}; host
\code{scipy.spatial.cKDTree}, device \code{lax.top\_k}), with weights
\begin{equation}
\label{eq:knn-weights}
\omega_n = \frac{1}{1 + (d_n / d_k)^2}, \qquad d_k = \max_n d_n ,
\end{equation}
and the design row of node $n$ at offset $\delta = \tilde p_n - \tilde p_q =
(\delta_x, \delta_v)$ (\code{BASIS = "quadratic\_full"}, \code{LocalFit.\_design},
device \code{ensemble\_device.\_design}):
\begin{equation}
\label{eq:design}
\phi(\delta) = \big(1,\;\; \delta_x,\;\; \delta_v,\;\; \{\delta_{v,a}\delta_{v,b}\}_{a \le b},\;\; \delta_x^2\big)
\in \R^{1 + D + 3 + 6 + D_x},
\end{equation}
so any field quadratic in $v$ (the adjoint of an energy-like reward) is represented
exactly.  (The host squares only the first spatial coordinate, $D_x = 1$; the device
squares every spatial coordinate, $D_x = D$; identical for $D = 1$.)  The weighted
least-squares coefficients solve
\begin{equation}
\label{eq:fit}
\Big(\textstyle\sum_n \omega_n \phi_n \phi_n^{\top} + \rho\, \bar d\, \operatorname{diag}(0, 1, \dots, 1)\Big)\beta
 = \sum_n \omega_n \phi_n \lambda_n,
\qquad \rho = \code{RIDGE\_REL} = 10^{-8},
\end{equation}
where $\bar d$ is the mean diagonal of the normal matrix (a \emph{relative} ridge on
every column but the intercept).  Then
\begin{equation}
\label{eq:field-value}
\Lam(x_q, v_q) = \beta_0, \qquad
\nabla_v \Lam = \beta_{1+D:\,4+D} / v_{\mathrm{scale}}, \qquad
\nabla_x \Lam = \beta_{1:\,1+D} / \Delta x .
\end{equation}
\code{ensemble\_device.field\_query} returns $(\Lam, \nabla_v \Lam)$; the $x$-slopes are
computed by the same solve but not returned.  The host
\code{LocalFit(gradient=True)} returns both slope blocks.  Only the value is consumed
by the shipped estimator; $\nabla_v \Lam$ and $\nabla_x \Lam$ are for the planned terms
of section~\ref{sec:shape}.  A query in a cell with no valid node returns~0.
The device \code{field\_query} docstring still says ``ridge \code{1e-12}''; the code is the
relative $10^{-8}$ of~\eqref{eq:fit} on both host and device, and the ridge is not
innocent (review section~5.12: the sequential gradient moves $3\times10^{-3}$ between
\code{RIDGE\_REL} $10^{-10}$ and $10^{-8}$ on the $N = 256$ inlet).

\section{Score terms and the parameter gradient}
\label{sec:scores}

Two estimators share the forward.  The \emph{trajectory} (DiCE) estimator of
\code{jax\_dsmc/adjoint.py} consumes the per-particle log-likelihood accumulator
\code{logh\_acc} that the forward fills; the \emph{ensemble} estimator of this note
consumes none of it.  This section lists both, so that ``which score is accumulated''
is unambiguous.

\subsection{What the forward accumulates in \code{logh\_acc} (trajectory route only)}
\label{sec:forward-scores}

\begin{enumerate}
\item NTC attempt count: $\log P(n_{\mathrm{att},c})$ of the stochastic rounding
  of~\eqref{eq:attempts} (\code{ntc.attempt\_count}), and per attempt
  $\log h = \log p_{ij}$ on acceptance, $\log(1 - p_{ij})$ on rejection (\code{vss.log\_h}).
\item Diffuse-wall accommodation Bernoulli (\code{boundary.diffuse\_from\_draws}):
  \begin{equation}
  \label{eq:acc-bernoulli}
  \log p_{\mathrm{acc}} = \ind[u \le \mathrm{acc}_f] \log \mathrm{acc}_f + \ind[u > \mathrm{acc}_f]\log(1 - \mathrm{acc}_f),
  \end{equation}
  while $T_f$ and $u_f$ are \emph{pathwise} there: the re-emitted velocity is the
  reparameterisation $v = \sqrt{k_B T_f/m}\,(z_t + \hat n\, r) + u_f$ of standard draws
  $(z, r)$ (\code{sampling.diffuse\_velocity\_from\_draws}), so \code{jax.grad} reaches
  $T_f$, $u_f$ through the velocity and $\mathrm{acc}_f$ through the score only.
\item The $\varepsilon$-hit interval log-probability (\code{wall\_adjoint.conditional\_time}):
  the realised event history restricts the Gaussian flight time
  $\tau \sim \mathcal N(\Delta t, \varepsilon^2)$ to an interval $[\tau_{\mathrm{lo}}, \tau_{\mathrm{hi}})$
  (last hit time, next hit time), and
  \begin{equation}
  \label{eq:eps-logp}
  \log p_\varepsilon = \log\Big[\operatorname{ndtr}\!\Big(\tfrac{\tau_{\mathrm{hi}} - \Delta t}{\varepsilon}\Big) - \operatorname{ndtr}\!\Big(\tfrac{\tau_{\mathrm{lo}} - \Delta t}{\varepsilon}\Big)\Big],
  \qquad \operatorname{ndtr} = \text{standard normal CDF},
  \end{equation}
  with the sample itself re-expressed as the value-preserving conditional quantile so
  that its tangent is the conditional one (\code{corrected = stop\_gradient(raw) + sampled - stop\_gradient(sampled)}).
  This is P3 section~4's random-time device; it is zero at $\varepsilon = 0$.
\item Inflow count (\code{inflow.inject}): $\log p_{\mathrm{count}} = \log\operatorname{frac}\bar N_f$
  if the rounding went up, $\log(1 - \operatorname{frac}\bar N_f)$ otherwise, the only path by
  which $n_\rho$ reaches the trajectory estimator.
\end{enumerate}

\subsection{What the ensemble estimator accumulates}
\label{sec:ensemble-scores}

All of the following weight a \emph{carried value at the post-move state} $\lambda^m$
(\code{Backward.lam}, \code{Births.values}, \code{Hits.values}); none of items~1--4
above is used, no $\varepsilon$-term appears, and no count score appears.

\paragraph{Diffuse walls: $\mathrm{acc}_f$, $T_f$, $u_f$ (all scored).}
The wall is a kernel and every recorded bounce $h$ of every slot in every step
(\code{WallEvents}, shape $(K_{\mathrm{bounce}}, N)$; newborn birth-step bounces included)
contributes
\begin{equation}
\label{eq:wall-score}
\frac{\partial J}{\partial \theta_w} = \sum_{h} \lambda^m_{i(h)}\; \partial_{\theta_w} \log b_{f(h)}\big(v^{\mathrm{out}}_h;\, \theta_w\big),
\qquad
b_f(v) =
\begin{cases}
\mathrm{acc}_f\, \dfrac{v_n\, e^{-m v_n^2/2k_BT_f}}{k_B T_f/m}\, \dfrac{e^{-m|v_t - u_f|^2/2k_BT_f}}{2\pi k_B T_f/m} & \text{diffuse branch},\\[2ex]
1 - \mathrm{acc}_f & \text{specular branch},
\end{cases}
\end{equation}
with the branch read from \code{WallEvents.diffuse} $= \ind[u \le \mathrm{acc}_f]$
(\code{examples/ensemble\_couette.py::log\_wall\_kernel}, assembled as
$\sum_h \code{values}_h \cdot \code{jax.grad}(\log b)$).  The weight is $\lambda$ at
the end of the move, not after this particular bounce, which is unbiased by the tower
property whatever the number of later bounces in the step (review section~5.6).
Note the difference from the trajectory route: there $T_f$ and $u_f$ are pathwise
through the reparameterised draw and only $\mathrm{acc}_f$ is scored; here the realised
re-emission is data and all three wall parameters enter through the kernel score.

\paragraph{Inflow sources: $n_\rho$, $T$, $u$ of a reservoir (density score + intensity).}
With $b_f$ the normalised inward flux density (\code{examples/ensemble\_inlet\_shock.py::log\_flux\_density})
\begin{equation}
\label{eq:flux-density}
b_f(v;\theta) = \frac{(v\cdot n_f)\, \exp\!\big(-m|v - u_f|^2/2k_B T_f\big)}{Z_f},
\qquad Z_f = \pi^{3/2} v_{s,f}^4\, F(s_f),\qquad v\cdot n_f > 0,
\end{equation}
\begin{equation}
\label{eq:source-gradient}
\frac{\partial J}{\partial \theta} =
\sum_{k=0}^{M-1}\sum_{f}\Big[
\sum_{b \in \mathcal B_{k,f}} \lambda^m_{b}\; \partial_\theta \log b_f(v_b;\theta)
\;+\;
\partial_\theta \bar N_f(\theta)\cdot \frac{1}{|\mathcal B_{k,f}|}\sum_{b \in \mathcal B_{k,f}} \lambda^m_{b}
\Big],
\end{equation}
where $\mathcal B_{k,f}$ are the slots born through face $f$ in step $k$ that survived
their birth flight (\code{Births}), $v_b$ is the velocity \emph{as sampled}
(\code{born\_v}), and $\bar N_f$ is~\eqref{eq:target-count} differentiated by
\code{jax.grad} through \code{inflow.target\_count}.  The first term is
\pref{ensemble\_inflow\_gradient\_left}/\pref{sampled\_left\_score\_gradient\_inflow}
with P3's $N_L/m_{L,k}$ weight equal to~1 (every birth carries its own value, no
interpolation); the second is the intensity derivative P3 states in words after
\pref{tangential\_temperature\_score\_inflow} and leaves out of its formulation
($N_L, N_R$ held fixed there).  It is exact for the stochastically rounded count
because $\E[N_f] = \bar N_f$ is smooth (review section~5.6, audit table).  For the inlet
shock every control acts through the Rankine--Hugoniot map, so
$\partial_\theta$ of both $b_f$ and $\bar N_f$ is taken by \code{jax.grad} through
\code{inlet\_shock\_reservoirs}.

\paragraph{Initial condition.}
The same two terms with $\lambda_{0,i}$ (\code{Adjoint.lambda0}, the value at $t_0$ on
the initially alive slots) against the initial Maxwellian's log-density and the
initial count intensity (\code{ensemble\_gradient} in \code{ensemble\_inlet\_shock.py}).

\paragraph{Collision-kernel parameters: \code{diam}, \code{omega}.}
\begin{equation}
\label{eq:kernel-gradient}
\frac{\partial J}{\partial \theta} = \sum_k \sum_{\text{pairs } p}\, W_p\, \Phi_p\, \partial_\theta \log q(v_{i_p}, v_{j_p}),
\qquad
W_p = \begin{cases} 1 & \text{realised pairs (sequential form)},\\ \tfrac12 w_{i\kappa} & \text{virtual pairs, \eqref{eq:kernel-term}},\end{cases}
\end{equation}
with $\partial_\theta \log q$ by \code{jax.grad} through \code{collide.species\_params}
(\code{examples/ensemble\_relaxation.py}).  Accepted pairs sample $q f f$ whatever the
majorant, so no attempt-count term appears for a kernel parameter.

\paragraph{Not covered.}
The VSS angular parameter $\alpha$ acts pathwise through \code{scatter} and would need
$\nabla_v\Lam$ at the post-collision states; \code{field\_query} returns it, no assembly
exists (review section~5.0).  Reparameterised (pathwise) source terms inside the
ensemble sweep were measured on the P3 model and rejected as biased and $1.8\times$
noisier (review section~5.9).

\section{Births and deaths}
\label{sec:births}

\paragraph{Births.}
A slot born in step $k$ is alive at $t_{k+1}$, collides in step $k$ like any other, and
so receives a value $\lambda^m_b$ from the collision pull-back like any other.  The
backward then does two things with it (\code{backward}, block ``2. sources''; device
values are read off \code{Backward.lam} by \code{from\_device}):
\begin{equation}
\label{eq:birth}
\code{Births.values}_b = \lambda^m_b = \lambda\big(x^m_{k,b}, v^m_{k,b}\big),
\qquad
\mu^{(k)}_b = 0 ,
\end{equation}
i.e.\ the value at the newborn's post-move state (its face position plus the
$\xi v_b$ first flight, after any birth-step bounces) is what the source score
of~\eqref{eq:source-gradient} multiplies, and the slot carries nothing further back
(it did not exist at $t_k$).  P3 says the same: an injected particle's adjoint
``contains its complete future effect'' and there is no score for $\xi$ because its
law is parameter-free.  A birth absorbed during its own first flight (\code{escaped})
is not alive at $t_{k+1}$ and is absent from \code{Births}, so it contributes~0 to the
score sum; the intensity mean in~\eqref{eq:source-gradient} is taken over the
surviving births only (the code's \code{mean(b.values[on])}), whereas the P3-model
route of review section~5.1 counts escaped births as zeros.  The two agree when no
birth escapes during its first flight; the difference is not measured on the
production cases.

\paragraph{Deaths.}
A particle absorbed during the move of step $k$ has $a^m_{k,i} = 0$, so
$\lambda^m_i = 0$ (\code{lam = where(alive, lam, 0)}) and by~\eqref{eq:move-pullback}
$\mu^{(k)}_i = 0$: the indicator of \pref{ensemble\_inflow\_transport\_adjoint}.
Slots that are dead throughout carry~0 and are never fit nodes (\code{cell\_nodes}
marks them invalid).  The device still \emph{queries} dead slots (wasted work on
capacity-padded runs, review section~5.12), which does not change any value.

\section{Validation summary}
\label{sec:validation}

Every number is copied from the review; ``FD'' is the same seeds' common-random-number
finite difference unless marked ``FD-only'' (independent seeds).  Sequential rows are
review section~5.7; virtual rows are section~5.12 (staged nodes, $K = 8$, device; $S = 2$
unless the row says otherwise, the relaxation rows from its pull-back scan table);
the last row is section~6.

Routes: \emph{src+int} = source density score plus intensity term~\eqref{eq:source-gradient};
\emph{wall} = diffuse-wall kernel score~\eqref{eq:wall-score}; \emph{kernel} = collision-kernel
score~\eqref{eq:kernel-gradient}; \emph{value} = delta-method terminal vector~\eqref{eq:delta-method}.
\begin{center}
\scriptsize
\begin{tabular}{llllll}
\hline
control & route & pull-back & ensemble $\pm$ SE & arbiter $\pm$ SE & $z$ \\
\hline
inlet \code{nrho}, $N{=}128$, $M{=}4$ & src+int & sequential & $1.375\mathrm{e}5 \pm 1.3\mathrm{e}3$ & FD $1.377\mathrm{e}5 \pm 2.5\mathrm{e}3$ & $-0.08$ \\
inlet \code{nrho}, $N{=}256$, $M{=}8$ & src+int & sequential & $1.361\mathrm{e}5 \pm 1.9\mathrm{e}3$ & FD $1.34\mathrm{e}5 \pm 2.3\mathrm{e}3$ & $0.71$ \\
inlet \code{temp}, $N{=}128$, $M{=}4$ & src+int & sequential & $1.345\mathrm{e}5 \pm 4.1\mathrm{e}3$ & FD $1.373\mathrm{e}5 \pm 1.1\mathrm{e}3$ & $-0.77$ \\
inlet \code{temp}, $N{=}256$, $M{=}8$ & src+int & sequential & $1.364\mathrm{e}5 \pm 3.6\mathrm{e}3$ & FD $1.374\mathrm{e}5 \pm 1.1\mathrm{e}3$ & $-0.32$ \\
inlet \code{temp}, $N{=}128$, $M{=}4$ & src+int & virtual & $1.302\mathrm{e}5 \pm 9.4\mathrm{e}3$ & FD $1.368\mathrm{e}5 \pm 2.6\mathrm{e}3$ & $-0.8$ \\
inlet \code{nrho}, $N{=}128$, $M{=}4$ & src+int & virtual & $1.377\mathrm{e}5 \pm 3.1\mathrm{e}3$ & FD $1.347\mathrm{e}5 \pm 6.0\mathrm{e}3$ & $+0.5$ \\
inlet \code{temp}, $N{=}2048$, $\Delta t{=}0.5$ & src+int & virtual & $1.361\mathrm{e}5 \pm 3.9\mathrm{e}3$ & FD $1.386\mathrm{e}5 \pm 1.5\mathrm{e}3$ & $-0.7$ \\
inlet \code{nrho}, $N{=}2048$, $\Delta t{=}0.5$ & src+int & virtual & $1.367\mathrm{e}5 \pm 1.6\mathrm{e}3$ & FD $1.375\mathrm{e}5 \pm 2.1\mathrm{e}3$ & $-0.4$ \\
Couette wall speed ($\varepsilon{=}0.005$) & wall & sequential & $0.5916 \pm 0.0097$ & FD $0.5999 \pm 0.0062$ & $-0.97$ \\
Couette wall speed ($\varepsilon{=}0$) & wall & sequential & $0.5865 \pm 0.014$ & FD $0.5933 \pm 0.0093$ & $-0.53$ \\
Couette $T_{\mathrm{left}}$ & wall & sequential & $0.1634 \pm 0.0088$ & FD $0.1804 \pm 0.012$ & $-1.2$ \\
Couette accommodation & wall & sequential & $0.2917 \pm 0.009$ & FD $0.2615 \pm 0.017$ & $1.7$ \\
Couette wall speed & wall & virtual & $0.615 \pm 0.030$ & FD $0.591 \pm 0.020$ & $+0.8$ \\
Couette $T_{\mathrm{left}}$ & wall & virtual & $0.169 \pm 0.018$ & FD $0.108 \pm 0.031$ & $+1.7$ \\
relaxation \code{diameter} & kernel & sequential & $-0.663 \pm 0.026$ & FD $-0.6723 \pm 0.012$ & $0.32$ \\
relaxation \code{omega} & kernel & sequential & $0.2824 \pm 0.01$ & FD $0.2607 \pm 0.014$ & $1.2$ \\
relaxation \code{diameter}, $S{=}4$ & kernel & virtual & $-0.6753 \pm 0.0070$ & FD-only $-0.6800 \pm 0.0034$ & ($+0.7\%$, $0.6\sigma$) \\
relaxation \code{diameter}, $S{=}8$ & kernel & virtual & $-0.6795 \pm 0.0071$ & FD-only $-0.6800 \pm 0.0034$ & ($+0.1\%$, $0.1\sigma$) \\
relaxation \code{omega}, $S{=}2$ & kernel & virtual & $0.2730 \pm 0.0039$ & FD-only $0.2706 \pm 0.0049$ & ($+0.9\%$, $0.4\sigma$) \\
RH shock $d(\rho_2/\rho_1)/dM_1$, 300 steps & value & virtual & $2.74 \pm 0.66$ & FD $2.54 \pm 0.17$ / $2.17 \pm 0.42$ & $0.3$ \\
\hline
\end{tabular}
\end{center}

\noindent
Further facts from the review used above: the ensemble's noise relative to the FD's
(sequential form, section~5.7) is $0.5$--$0.8\times$ for reservoir density, $0.7\times$ wall
temperature, $0.7\times$ \code{omega}, $1.5\times$ wall speed, $2.1\times$ diameter,
$3.3$--$3.8\times$ reservoir temperature; the virtual form's per-seed sd is $7\times$ below the
sequential form's at $0.16$ collisions/particle/step and $37$--$56\times$ at $0.8$
(section~5.12); the accommodation FD is a step function of a decision parameter and
cannot arbitrate below $\sim 3\%$ at 384 seeds (section~5.6); the RH stationary slope
$1.96$ is a reference only; the pathwise VJP of the same RH functional had the wrong
sign ($-1.9\mathrm{e}{-3}$ vs FD $+1.8\mathrm{e}{-3}$, \code{examples/shock/gradients.py}).
Host and device agree to $10^{-8}$ on identical records for the sequential form and
statistically for the virtual form (\code{tests/test\_ensemble\_device.py}); the
virtual pull-back carries a quadratic invariant to $10^{-4}$ at the shipped ridge and
$10^{-12}$ at zero ridge (\code{test\_virtual\_pullback\_carries\_an\_invariant\_exactly}).
The shape oracle of section~\ref{sec:shape} is in \code{results/shape\_planar\_oracle.json}.

\section{The shape adjoint (implemented 2026-09-18: \code{jax\_dsmc/shape.py})}
\label{sec:shape-built}

The geometry $\theta$ (the centre and radius of an analytic body, or the vertex
coordinates of a surface; normals and areas are rebuilt from it, the cut volumes and
cell tables are not) enters the forward only through the move of each step, which
under the RNG contract is a pure map
\begin{equation}
\label{eq:wall-map}
(x^{\pth}_k, v^{\pth}_k, \log p_{\varepsilon,k}, T_k, \mathcal{E}_k)
 = \mathcal{W}_\theta\big(x^{\pre}_k, v^{\pre}_k, \kappa_k\big),
\end{equation}
of the previous post-collision state and the step's key $\kappa_k$ (\code{key\_for(base,
k, STAGE\_MOVE)}); $T_k$ are the step's surface tallies and $\mathcal{E}_k$ its traced
events (hit point $x^h$, incident velocity $v^{\mathrm{in}}$, normal $n$, event time
$t^h$ of every recorded hit, \code{advance(trace=True)}).  \code{shape.replay} re-runs
$\mathcal{W}_\theta$ on the captured record and checks it reproduces the record
(relative error $<10^{-12}$: the same arithmetic in another compilation);
\code{shape.gradient} takes \code{jax.vjp} of~\eqref{eq:wall-map} with respect to
$\theta$ at every captured step and sums
\begin{equation}
\label{eq:shape-gradient}
\frac{\dd J}{\dd\theta} \approx \sum_k \Big[
 \underbrace{\nabla_x\lambda^m_k \cdot \partial_\theta x^{\pth}_k
 + \nabla_v\lambda^m_k \cdot \partial_\theta v^{\pth}_k}_{\text{pathwise } x, v}
 + \underbrace{\partial_T J \cdot \partial_\theta T_k}_{\text{tally}}
 + \underbrace{w_k\, \partial_\theta \log p_{\varepsilon,k}
 + \sum_j w_{k,j}\, \nabla_n \log q(v^{\mathrm{out}}_{k,j} \mid n_{k,j})\cdot \partial_\theta n_{k,j}}_{\text{score}}
 + \underbrace{\sum_j c_{k,j}\, \partial_\theta t^h_{k,j}}_{\text{timing}}
\Big],
\end{equation}
every $\partial_\theta(\cdot)$ being what the VJP of $\mathcal{W}_\theta$ delivers.

\paragraph{Cotangents.}
$\nabla_x\lambda^m_k$, $\nabla_v\lambda^m_k$ are the slopes of the per-cell
fit~\eqref{eq:field-value} at each survivor's own post-move state
(\code{field\_query} now returns \code{(value, grad\_x, grad\_v)};
\code{node\_gradients}).  The fit is \emph{one-sided}: only nodes with
$v\cdot n > 0$ for the hit particle's last normal $n$ are used, because next to a wall
the post-move population splits into particles about to hit (value $\approx$ the hit's
reward) and particles that just left, an $O(r)$ contrast a polynomial in $v$ cannot
represent and that is correlated with the distance to the wall (the approaching side
is depleted within one flight of it) -- the two-sided $x$-slope read that population
mix, and the plane's shear gradient came out with the wrong sign ($-0.012 \pm 0.005$
against FD $+0.017 \pm 0.003$).  $w_k$ is~\eqref{eq:hit-weight} summed over the step's
hits; $\partial_T J$ is the objective's weight on the tally (a windowed drag:
$f_{\mathrm{num}}/(W\,\Delta t)$ on the element rows' $x$-momentum).  The diffuse
re-emission velocity is \code{stop\_gradient} inside the map: a kernel's parameters
\emph{and} its normal enter only through the score of the frame-independent density
\begin{equation}
\label{eq:diffuse-density}
q(v\mid n) = \frac{v_n}{s^2}\,e^{-v_n^2/2s^2}\;\frac{e^{-|v_t - u|^2/2s^2}}{2\pi s^2},
\qquad s^2 = kT_w/m,\; v_n = v\cdot n,\; v_t = v - v_n n
\end{equation}
(\code{sampling.log\_diffuse\_density}); the specular branch and the hit point are
pathwise (review section~5.9).

\paragraph{The hit-timing band term.}
The VJP differentiates a \emph{fixed} hit set.  A reward that changes at a hit --
a velocity-dependent per-step reward, or a per-hit reward at the window's edges --
also depends on which step a hit falls in, and moving the surface moves hits across
the step boundary.  With $\Delta J_{k,j}$ the change of the objective when hit $j$
moves from the end of step $k$ to the start of step $k+1$
($\phi_k(v^{\mathrm{out}}) - \phi_k(v^{\mathrm{in}})$ for a state reward; $+r$ on the
last window step and $-r$ on the step before the window for a per-hit reward; the
value after either outcome is equal in distribution, so nothing else jumps), the
boundary term is the density of hit times at the step end times $-\Delta J$,
estimated from the hits within a band $\beta\,\Delta t$ of it:
\begin{equation}
\label{eq:timing}
c_{k,j} = -\frac{\Delta J_{k,j}}{\beta\,\Delta t}\,\ind\big[t^h_{k,j} > (1-\beta)\Delta t\big],
\end{equation}
placed on the traced event time so that $\partial_\theta t^h$ is exact for any body.
It is the value-preserving analogue of the $\varepsilon$ hit-timing score
of~\eqref{eq:eps-logp} with no forward regularisation.  On the plane's windowed $v_y^2$
the pathwise term alone was $+0.82 \pm 0.23$ against FD $+0.05 \pm 0.15$; the band term
is $-0.43 \pm 0.02$ ($\beta = 0.2$ and $0.4$ agree to 1\,\%).  A per-hit reward's term
lives at the two window edges and has zero mean on a steady window (the cancelling
edge effect of gate G3a), so \code{examples.ensemble\_mesh} reports it and leaves it
out of the estimate.

\paragraph{Estimating $\nabla_x\lambda$: the one-sided, pooled and along-flight forms.}
The cotangent $\nabla_x\lambda^m_k$ is the weak link of~\eqref{eq:shape-gradient}: next
to a surface the post-move population splits into particles about to hit (value
$\approx$ the hit's reward) and particles that just left, an $O(r)$ contrast a
polynomial in $v$ cannot represent, and the per-hit signal of a translated wall is a
window-end quantity far below the fit's slope noise.  Four estimators are implemented
and compared against replica-exact gradients (review section~7.6, item~5):
\begin{enumerate}
\item \emph{one-sided}: the fit~\eqref{eq:field-value} over the nodes with
  $v\cdot n > 0$ for the query's last hit normal $n$ (\code{node\_gradients(direction=)});
  the general rule for any query uses the nearest surface within one cell width
  (\code{mesh.geometry.nearest\_normal}, \code{backward\_virtual(normals=)});
\item \emph{pooled}: the $k$ nearest one-sided nodes drawn from the $3^D$ block of cells
  around the query's cell (\code{pool=1}, $k = 96$; \code{ensemble\_device.pooled\_tables});
\item \emph{along-flight}: collisionless free flight makes $V(x - s v, v)$ the value of
  the same particle delayed by $s$, so
  \begin{equation}
  \label{eq:along-flight}
  v\cdot\nabla_x V(x, v) = -T(x, v),\qquad
  T(z) = \E\Big[\sum_{k' > k}\sum_j \tau_{k',j}\,\Big|\,z_k = z\Big],\qquad
  \tau_{k,j} = c_{k,j}\ \text{of~\eqref{eq:timing}},
  \end{equation}
  $T$ being the value of the \emph{timing rewards} -- the change of the objective per
  unit delay of the trajectory -- carried by a second value-form backward with
  \code{hit\_reward} $= \tau$ (\code{Replay.timing\_reward}).  The cotangent then
  replaces the fitted component along $v$ by the carried value,
  $\nabla_x\lambda \leftarrow \nabla_x\lambda - \frac{(v\cdot\nabla_x\lambda) + T}{|v|^2}\,v$
  (\code{shape.along\_flight}), and keeps the fitted perpendicular part.  With
  collisions the identity acquires the collision operator's term ($v\cdot\nabla_x V =
  -\text{rate} - Q^*V + \partial_t V$); it is neglected here (the cases run at
  $\le 0.3\,\%$ collisions per particle-step);
\item \emph{pooled $+$ along-flight}: (2) for the perpendicular part, (3) along $v$.
\end{enumerate}
Since $\tau$ is nonzero only at band hits, the carried $T$ is a sparse sample of its
expectation: unbiased, but the per-hit reward at the window's two edges is what a per-hit
objective's $T$ consists of, while a velocity-dependent per-step reward has band hits at
every step.  Measured (review section~7.6, item~5; replica-exact gradients on the plane
and the production runs): the along-flight form removes the fitted estimator's bias on
the windowed $v_y^2$ ($-0.372 \pm 0.067 \to -0.188 \pm 0.040$ against the truncated FD
$-0.149 \pm 0.037$) at half the variance; pooling lowers the per-seed sd 4--7$\times$ on
the circle radius but on its own keeps a bias ($-0.20 \pm 0.11$ on the $v_y^2$ x-part);
the combination (4) is the default of the gradient runners.  What none of them resolves
is the x/v-part of a translated flat wall: on the shear probe the exact value is
$+0.0032 \pm 0.0006$ and the pooled / along-flight forms read $0$ (pooling smooths a
slope a tenth of the per-particle noise to zero; the realised $T$ is a sparse sum of
window-edge band hits whose sample SE understates its uncertainty), while the one-sided
fit is consistent at $3\times$ the noise.  Radius-type controls are FD-consistent to a
few per cent because their fixed-hit-set tally term is exact; translation-type
controls of flat walls are not resolved by the fitted or the value-based forms here.
\emph{Decision (open issue):} shape controls are therefore restricted to normal
displacements of curved surfaces, which is sufficient for the first optimisation case
(a circle deformed toward an airfoil-like shape, and the sphere analogue); tangential
translations of flat walls are not supported as controls.

\paragraph{Stratifying the fit by the hit-ahead class.}
The one-sided rule separates the two populations exactly at a flat wall and not at a
curved one: a particle leaving a cap's tilted tangent plane can still be heading for the
flat face below it, so the pool of a reflected query mixes in nodes that carry their
next hit's reward, and the fitted $\partial\lambda/\partial v_x$ reads the mixture's
$d(2 m v_x)/dv$ times the hit probability.  On the specular capsule (review
section~7.9) this is a $-27\,\%$ bias of the total against the CRN FD, entirely in the
$\nabla_x\lambda$, $\nabla_v\lambda$ parts, whose true value is zero: for a convex body
in a collisionless flow a reflected particle never re-hits, and the own-reward (tally)
term alone reproduces the FD seed by seed.  The record holds every particle's future,
so the class is exact: the nodes of a query are the particles of its own class,
\begin{equation}
\label{eq:hit-ahead}
\kappa_k^m = \mathbf 1\{\exists\, k' > k:\ \text{particle } m \text{ hits the body at
step } k'\},
\qquad \text{nodes of } (k, m) = \{\,m' : \kappa_k^{m'} = \kappa_k^m\,\}
\end{equation}
(\code{shape.hit\_ahead}, \code{node\_gradients(klass=)}).  For the collisionless
convex bodies the post-hit class has $\lambda \equiv 0$ and the $\nabla\lambda$ parts
vanish; in a collisional gas the class-conditional fit is the smooth part of the value
gradient, the class boundary (a hit-set change) being what the silhouette and timing
terms carry.  Specular capsule, four CRN seeds: $0.964 \pm 0.037 \to 1.356 \pm 0.011$
against the FD $1.35 \times 10^{-20}$ ($z = -11 \to +0.08$), the x/v parts
$-3.9 \to +0.05 \pm 0.6 \times 10^{-21}$.

\emph{Representation.} The shape is optimised as a $C^1$ level set
$\phi(x;\,c) = \sum_j c_j B_j(x)$ (uniform cubic B-splines on a grid finer than the
DSMC grid; body $\phi<0$, normal $n=\nabla\phi/|\nabla\phi|$ into the flow), not as
polygon vertices.  For a flat-normal polygon the Hadamard derivative carries a delta at
every vertex (the normal jumps there); under uniform scaling that corner term is the
whole $\mathrm dF/\mathrm dr$.  On a $C^1$ surface the hit normal varies smoothly with
the shape and the pathwise derivative of the replayed wall map is complete for specular
walls; a $C^0$ interpolant would merely move the jumps to the level-set cell faces.  The
hit time is the root of $\phi(x+tv)=0$, found by bracketing and bisection and finished by
one Newton step with the bracket held constant, so that
$\partial t/\partial c_j = -B_j(x_{\rm hit})/(v\cdot\nabla\phi)$ (the implicit-function
derivative; the bisection alone has none).  The only non-pathwise term that survives is
the silhouette term of a diffuse body (a near miss becoming a grazing hit), on the set
$\min_s \phi(x+s\hat v)\in(0,\varepsilon)$; it vanishes for specular walls.

\paragraph{The silhouette band term.}
Let $\phi_{\min}(\omega) = \min_s \phi(x + s\hat v)$ be the closest approach of a
particle's undeflected flight to the surface ($<0$: it hits, with that chord depth;
$>0$: it misses by that distance).  A displacement of the surface changes the outcome of
the particles with $|\phi_{\min}| < |\delta\phi|$; the change of the objective is
$V_{\rm hit} - V_{\rm miss}$, the value of the realised outcome being the particle's
reward-to-go and that of the virtual one (a straight flight for a hit, a re-emission from
the closest point under the wall model for a miss) the fitted adjoint at the virtual
post-move state plus the virtual hit's reward.  With $\rho(0)$ the density of
$\phi_{\min}$ at zero,
\[
  \frac{\mathrm dJ}{\mathrm d\theta}\Big|_{\rm sil}
  = -\,\mathbb E\Big[\rho(0)\,(V_{\rm hit}-V_{\rm miss})\,
      \frac{\partial \phi_{\min}}{\partial\theta}\Big]
  \approx -\frac{1}{2\varepsilon_b}\sum_{|\phi_{\min,i}|<\varepsilon_b}
      (V_{{\rm hit},i}-V_{{\rm miss},i})\,B_j(x^*_i)\,\frac{\partial c_j}{\partial\theta},
\]
the band on both sides of the surface (window $2\varepsilon_b$) estimating $\rho(0)$,
and $\partial\phi_{\min}/\partial c_j = B_j(x^*)$ by the envelope theorem at the frozen
closest point $x^*$.  Every passing is counted once: a miss in the step whose move
contains $x^*$ (strictly inside -- a particle starting there has just left the
surface), a hit in the step it hits, with $\phi_{\min}$ taken along the line within
the search reach.  For a specular wall $V_{\rm hit}\to V_{\rm miss}$ at grazing
incidence and the term is $O(\varepsilon_b)$; for a diffuse wall it is the $O(1)$ term of
the silhouette, and on constructed rays with $\lambda = 0$ it reproduces (new hits per
unit displacement) $\times$ (absorbed momentum) to $0.04\,\%$ at $\varepsilon_b = 0.02$.

\paragraph{Grazing regularisation.}
The pathwise derivative of a hit carries $\partial t/\partial\theta = -\,\partial_\theta\phi
/ (v\cdot\nabla\phi)$; over the ray measure $\mathrm db$ its mean converges but its
variance diverges logarithmically at grazing incidence (one hit at $|v\cdot n|/|v| =
3\cdot10^{-3}$ was $15\,\%$ of a window's tally sum).  The denominator is evaluated as
$\max(|v\cdot\nabla\phi|,\ \epsilon\,|v|\,|\nabla\phi|)$ in the tangent only (the primal
hit is exact): the per-hit weight is bounded by $1/\epsilon$ and the bias is
$O(\epsilon)$ -- measured on the circle's radius tally term over 8 CRN seeds, the seed
sd falls from $8.4\,\%$ to $4.2\,\%$ ($\epsilon = 0.03$), $2.9\,\%$ ($0.1$), $2.1\,\%$
($0.3$) with paired biases $+0.5\pm1.8$, $-1.2\pm2.3$, $-7.4\pm2.4\,\%$.
The diffuse normal score has the same tail on the emission side: $\partial_n \log q =
v_{\mathrm{out}}/v_n - v_n v_{\mathrm{out}}/s^2 + \dots$ with $v_n$ Rayleigh, so
$\mathbb E[(v_t/v_n)^2] = \int \dd v_n\,(v_n/s^2)\,|v_t|^2/v_n^2$ diverges
logarithmically at $v_n \to 0$ (the flux-weighted cosine law puts mass $\epsilon^2$
below $v_n/|v| = \epsilon$).  The derivative of $\log v_n$ is evaluated at
$\max(v_n,\ \epsilon |v_{\mathrm{out}}|)$ with the same $\epsilon$ (primal exact); the
bias of the clipped part is $\int_0^{\epsilon s} (v/s^2)(1/v - 1/\epsilon s)\,\dd v
\cdot |v_t| = \epsilon |v_t| / 2s$, i.e.\ about $\epsilon/2$ of the $v_t/v_n$ part.

\paragraph{Hybrid form of the diffuse normal term.}
With the per-hit reward $r(v_{\mathrm{out}})$ exact, the normal dependence of a diffuse
hit splits as
\[
\frac{\dd}{\dd n}\,\mathbb E\big[r(v_{\mathrm{out}}) + V(v_{\mathrm{out}})\big]
= \mathbb E\Big[\frac{\partial r}{\partial v_{\mathrm{out}}}\cdot
\frac{\partial v_{\mathrm{out}}}{\partial n}\Big]
+ \mathbb E\big[V(v_{\mathrm{out}})\,\partial_n \log q(v_{\mathrm{out}}\mid n)\big],
\]
the first term reparameterised through $v_{\mathrm{out}} = \sigma\,(r\,n + z_1 t_1(n)
+ z_2 t_2(n)) + u$ at fixed draws (weights bounded by $m|v|\,|\partial n/\partial\theta|$),
the second the score weighted by the value-to-go \emph{after} the hit, which needs no
fitted derivative of $\lambda$.  Because the tangential law is isotropic any smooth
frame field $t_{1,2}(n)$ gives the same pushforward, and the derivative is taken in the
frame parallel-transported from the primal one, $t_1(n) = (a - (a\cdot n)n)/|\cdot|$
with $a$ the frozen primal tangent, whose derivative $-(a\cdot\delta n)\,n$ is bounded;
the incident-velocity frame of the sampler has $\partial t_1/\partial n \sim 1/|v_{t,\mathrm{in}}|$
at normal incidence, a log-divergent variance for nothing.

\paragraph{What the terms are on the analytic circle.}
For a specular cylinder of radius $r$ in free-molecular flow the fixed-hit-set part of
$dF/dr$ is the whole of $F/r$: with impact parameter $b$, $\cos\varphi = \sqrt{1-b^2/r^2}$
and $F \propto \int_{-r}^{r}\cos^2\varphi\,\dd b = 4r/3$, the derivative at fixed rays
$\int 2b^2/r^3\,\dd b = 4/3$ equals $F/r$ while the new grazing rays contribute
$\cos^2\varphi = 0$.  Measured (SPARTA circle at $f_{\mathrm{num}}\times 4$, 700 warm-up
steps, 200-step window): the tally term is $1.27 \pm 0.02 \times 10^{-20}$ against
$F/r = 1.24\times 10^{-20}$; the flow-field feedback of the blockage ($\mathrm{FD} - F/r
\approx 0.15\times 10^{-20}$) is what the $\nabla\lambda$ terms carry, and the $x$-part
carries most of the estimator's noise (see the review, section~7.6).

\section{The shape extension as first planned (superseded by section~\ref{sec:shape-built})}
\label{sec:shape}

This paragraph is the plan review section~6 assigned to a wall-position control
$\theta$ before the terms were built; kept for the code comments that point at it.
(i)~\emph{Hit set}: the $\varepsilon$-hit interval log-probability
$\log p_\varepsilon$ of~\eqref{eq:eps-logp}, differentiated with respect to the wall
position instead of being ignored, times the carried value at the post-move state,
$\sum_i (\lambda^m_i + r^{\mathrm{hit}}_{k,i}/\Nn)\, \partial_\theta \log p_{\varepsilon,i}$, the per-hit
reward $r_k$ of section~\ref{sec:objective} riding along; the collisionless planar
oracle (\code{examples/shape\_planar\_oracle.py}) recovered $0.588 \pm 0.008$ against the
analytic $0.5917$ with this value-weighted form, where the frozen-hit pathwise
gradient gave~$0$.  (ii)~\emph{Reflected position}: a replayed wall-map VJP with the
cotangents $(\nabla_x \lambda, \nabla_v \lambda)$ from~\eqref{eq:field-value} at the post-hit
state, $\nabla_x\lambda\cdot \partial x'/\partial\theta$ with $\partial x'/\partial\theta = 2$
for a specular plane, requiring \code{field\_query} to return the $x$-slopes it already
computes.  (iii)~An \emph{edge-band score} for what the VJP cannot see: on a
flat-normal polygon every element's normal is constant under a scaling of the body, so
the whole $dF/dr$ sits in \emph{which} element a ray meets ($O(1/E)$ jumps at $O(E)$
vertices), and a diffuse body has an $O(1)$ silhouette term (a grazing hit re-emits
thermally; for a specular body it vanishes since $v_{\mathrm{out}} \to v_{\mathrm{in}}$ at
grazing incidence).  For hits within a band $\varepsilon_b$ of an element edge and
near-misses within $\varepsilon_b$ in impact parameter, the alternative outcome is
evaluated through the fit and the term is $(\Lambda_{\mathrm{alt}} - \Lambda_{\mathrm{actual}})\,
\partial_\theta(\text{edge position})/\varepsilon_b$, value-preserving like the
conditional time and reported separately.  (iv)~The cut-volume route: the NTC weight
$2 n_{\mathrm{att},c}/N_c$ of a cut cell depends on $V_c(\theta)$ through~\eqref{eq:attempts},
entering as $-\partial_\theta \log V_c$ times the substep's four-term sum
in~\eqref{eq:virtual} (the \code{pair\_score} slot); built only if the Kn~$\approx 0.5$
probe shows it matters.
There is \emph{no} cell-membership term: the cells are fixed (the body is embedded in
the Cartesian grid, SPARTA-style), so a shape change never moves a cell face.
Arbiters, in order: an embedded specular plane in a fixed box (position as control)
against ensemble FD over independent keys, $|z| \le 2$ at two $\Delta t$; then the analytic
circle's radius (free-molecular $dF/dr = F/r$ for a cylinder, $2F/r$ for a sphere, in an
unbounded stream); then the 50-gon through the edge-band score.

\end{document}
